\centerline{\bf \Huge ДЗ. Задачи на логику I}
\centerline{\bf \Huge }

\begin{flushright}
        \scriptsize{\textbf{Сайгаки пьют калмыцкий чай?!}}
\end{flushright}

\problem{\textbf{1}} На острове рыцарей и лжецов живут рыцари, которые всегда
говорят правду, и лжецы, которые всегда лгут. Однажды пятерых жителей этого острова по
очереди спросили, сколько среди них лжецов.

— Три, — ответил первый.

— Три, — ответил второй.

— Один, — ответил третий.

— Поверьте им, они все рыцари, — сказал четвёртый (про 1-3).

— Не верьте им, они все лжецы! — сказал пятый (про 1-4).

Сколько рыцарей было на самом деле?


\problem{\textbf{2}} Шесть сайгаков сидят за круглым столом. Известно, что ровно два
сайгака всегда говорят правду, и они сидят рядом. Кроме этого, ровно два сайгака всегда врут, и
они тоже сидят рядом. Оставшиеся два сайгака могут как врать, так и говорить правду, и они не
сидят рядом. Начальник ЭлМШ ходит вокруг стола и спрашивает сайгаков, где они спрятали калмыцкий чай.

• Первый сайгак сказал, что в Элисте.

• Второй сказал — в Ики-Бурульском районе.

• Третий сказал — в Ики-Бурульском районе, а именно в посёлке Оргакин (данный посёлок находится в Ики-Бурульском районе).

• Четвёртый сказал — в Ики-Бурульском районе, а именно в посёлке Бага-Бурул (данный посёлок находится в Ики-Бурульском районе).

• Пятый сказал — в Троицком.

Где сайгаки спрятали калмыцкий чай? Ответ нужно обосновать.


